% A coefficient alphabet independent of the represented set.
\section{Fixing the coefficient digit base at four}\label{sec:fixed-four}

The coefficient digit base can be fixed at $Z=4$ by expanding the
represented equations into elementary arithmetic gates. Their number
affects the fixed index, but not the number of unknowns in the final
system: the additional coordinates are digits of the existing code.
The exponent $L$ becomes a fixed index parameter in place of $Z$.
This identifies the already supplied $\theta=B-4$ with a separately
computed register of Theorem~\ref{thm:best102}, saving one addition.

\subsection{Primitive rows with a fixed coefficient alphabet}

Start with a finite polynomial system representing the desired set
over nonnegative witnesses and a positive input $x$. Split each equation
into an equality of polynomials with nonnegative coefficients, and
evaluate those polynomials by finite circuits of additions,
multiplications, copies, and constants zero and one. Every gate has a
fresh nonnegative output coordinate. Arbitrary integer coefficients are
implemented by finite arithmetic from one. Thus this construction
places no bound on the coefficients of the original system.

Replace repeated operands of an addition or multiplication by distinct
copied coordinates. With a separate homogenizing coordinate $\delta$,
the quadratic residuals have the following forms:
\[
\begin{array}{c|l}
\text{gate or constraint}&\text{homogeneous residual }F\\ \hline
\text{multiplication}&XY-W\delta\\
\text{addition}&(X+Y-W)\delta\\
\text{copy or equality}&(X-Y)\delta\\
\text{zero}&X\delta\\
\text{unit}&X\delta-\delta^2.
\end{array}
\]
Here $X,Y$ may include the input $x$, but are distinct from $\delta$;
outputs are fresh, and trivial identical equalities are omitted.
At $\delta=1$, these rows express the original circuit exactly.
Every solution of the original system extends by its nonnegative gate
values. Adjoin a fresh coordinate $u$ and the guard
$F_{\mathrm g}=u\delta-x^2$.

For every ordinary row, including the guard, use the target
\begin{equation}\label{eq:four-targets}
 G_j=2F_j+\delta^2,
 \qquad G_{\mathrm{special}}=\delta^2.
\end{equation}
For a homogeneous quadratic monomial $I$, let $c_I=1$ for a square
and $c_I=2$ for a product of distinct coordinates. Every coefficient
of every $G_j$, divided by $c_I$, belongs to
\begin{equation}\label{eq:four-coefficient-alphabet}
 \{-2,-1,0,1\}.
\end{equation}
Indeed, ordinary cross terms give $\pm1$, and the added $\delta^2$
gives $1$. The unit row gives $2X\delta-\delta^2$, with divided
coefficients $1,-1$. The guard gives
$2u\delta-2x^2+\delta^2$, with divided coefficients $1,-2,1$.
The special target has coefficient $1$. Distinct copied operands
prevent a repeated positive square or coincident cross terms from
enlarging this alphabet. The displayed orientation of the unit row
is part of the construction.

\subsection{Support and admissible indices of arbitrary finite size}

Let $m$ be the number of coordinates other than the input, before
adding dummy row coordinates. It includes $\delta$, all circuit
coordinates, and $u$. Give $x$ weight zero, and give the remaining
coordinates weights
\[
 v_i=3^i\quad(0\le i<m),\qquad v_0=1\text{ for }\delta,
 \qquad M=3^{m-1}.
\]
All quadratic monomial weights are distinct. The weights are zero,
a single $3^i$, or a sum of two such powers; their base-three digits
recover the unordered pair, including repeated coordinates and the
weight-zero input.

Let $s$ be the total number of targets, including the guard and the
special target. Define
\begin{align*}
 d_0&=4M+1,& t_0&=(s+1)d_0+2M,\\
 t_j&=t_0+jd_0\quad(0\le j<s),&
 t_*&=2sd_0+2M,& K&=t_*+1.
\end{align*}
Choose a power of two $L>3K+2$. The coefficient intervals
$[t_j-2M,t_j]$ are disjoint and lie above $M$. Moreover
\[
 2t_0-2M>t_*.
\]
Include a dummy coordinate at each row position $t_j$. This last
inequality makes every contribution involving a dummy lie above all
targets, so the dummies cannot alter a tested residual.

Write $a_{j,I}$ for the coefficient of monomial $I$ in $G_j$, and set
\begin{align*}
 \mathcal D(B)&=\sum_{j,I}\frac{a_{j,I}}{c_I}B^{t_j-w(I)},\\
 \ell_0(B)&=\sum_{i=0}^{m-1}B^{v_i}+\sum_{j=0}^{s-1}B^{t_j},\\
 e_0(B)&=2\sum_{j=0}^{K-1}B^j+\mathcal D(B),&
 V&=\ell_0(4)+e_0(4)4^L.
\end{align*}
Monomial uniqueness, disjoint row intervals, and the dummy inequality
give
\begin{equation}\label{eq:four-coefficient-identity}
 [B^{t_j}]\,\mathcal D(B)\mathcal C(B)^2=G_j.
\end{equation}
They also show that the coefficients of $\mathcal D$ lie in
\eqref{eq:four-coefficient-alphabet}. Consequently all digits of $e_0$
belong to $\{0,1,2,3\}$. Zero digits are permitted. Since $\mathcal D$
has no constant term, the unit digit of $e_0$ is exactly $2$; in
particular $e_0$ and $V$ are positive.

Put $c_*=m+s+1$, the number of code coordinates including the input.
Choose a power of two
\begin{equation}\label{eq:four-index-bound}
 H>\max\{2\,4^{2L+1},\ 4^{t_*+3}\,4c_*^2,\ 3L,\ 16\}.
\end{equation}
An admissible index is a triple $(V,H,L)$ obtained by this construction
from a finite representing system. All choices, including the circuit
size, are made independently of the queried input $x$. The supplied
positive parameters are exactly $(x,V,H,L)$. The variable number of
decoded coordinates introduces no new unknown in the final system.

\subsection{The fixed-four system and its preliminary bounds}

In the system of Theorem~\ref{thm:best102}, set $Z=4$ and replace the
former fixed exponent by the index parameter $L$. With the same charged
abbreviations and positive witnesses, the coding equations become
\begin{align*}
 B&=Hb^2,& b&=x+\beta,& \lambda(B-1)&=q^2-1,& \theta+4&=B,\\
 S_2&=\ell+eq=V+t\theta,& S_2+\alpha&=q^2,&
 \Omega&=4\lambda-2e>0,&\mathcal C&=x+g,\\
 S_3&=B\lambda q-\Omega\mathcal C^2,& \sigma&=S_3>0,&n&=q^8.
\end{align*}
The packing is
\begin{align*}
 S&=g+q^2(S_2+q^2S_3),\\
 T+1&=q^2(1+\theta\lambda)-(b-1)\ell+\theta\ell q^4,\\
 r&=S(n^2-n)+(T+1)(n^2-1).
\end{align*}
All common-witness Pell equations of Section~\ref{sec:common-witness}
are retained. In the second exponential congruence the expression
$a-(B-4)$ is now $a-\theta$.

Before decoding, $S_2<q^2$ gives $e<q$ and $\ell<q^2$.
The geometric equation and $H>16$ give
$B\le q^2$, $q>4b$, and $3L\le B\le n$.
Since $\Omega\ge1$ and $S_3>0$,
\[
 \mathcal C^2<B\lambda q<3q^3<q^4.
\]
The first block borrows fewer than $b$ units from the middle block;
the latter remains positive because $\theta=B-4>b$.
The width-$(2,2,4)$ argument of Section~\ref{sec:borrow-mask} therefore
gives
\[
 0<S<n,\qquad 0<T+1\le n,\qquad n\le r\le2n^3-2n^2.
\]
These are the hypotheses of the common-witness Pell argument. Its
exponent comparisons use $L\ge2$ and $3L\le B\le n\le r$, which
hold for every admissible index here. It follows that $b,B,q$ are
powers of two, $q=B^L$, and the required central-binomial divisibility
holds. No bound on the number of compiled rows is used in that argument.

\subsection{Canonical codes, zero digits, and the target tests}

Now $\lambda>q$, so $\Omega>2\lambda$. Positivity of $S_3$ gives
$\mathcal C^2<Bq/2<q^2$, hence $g<\mathcal C<q$.
Let $d=\lfloor(b-1)\ell/q^2\rfloor<b$. Subtracting $d$ from
$(B-4)\lambda$ changes only its unit digit. The middle binary mask
therefore bounds all nonunit digits of $S_2$ by three and its unit
digit by $3+d$. Thus
\[
 \ell\le S_2<\frac{4q^2}{B-1}+b.
\]
Using $B=Hb^2$ and $H>16$ shows $(b-1)\ell<q^2$, so $d=0$.
The canonical evaluation lemma applies to
$\ell_0(B)+e_0(B)B^L$, whose digits are below four and whose degree
is below $2L$. The first term of \eqref{eq:four-index-bound} supplies
the required evaluation bound. It yields
\[
 S_2=\ell_0(B)+e_0(B)q,\qquad
 \ell=\ell_0(B)+a_{\mathrm{al}}q,\qquad e=e_0(B)-a_{\mathrm{al}},
 \qquad0\le a_{\mathrm{al}}<e_0(B).
\]
The quantity $a_{\mathrm{al}}$ is the temporary quotient alias.

The first mask, with $g<q$, decodes the low positions specified by
$\ell_0$. Every coordinate is below $b$, and the unit digit of
$\mathcal C$ is $x$. Put
\[
 A_{\mathcal C}=\mathcal C(1)^2,
 \qquad J_{\mathcal C}=4A_{\mathcal C}.
\]
There are $c_*$ coordinates, so
$\mathcal C(1)<c_*b$ and $J_{\mathcal C}<B/4$.
Also $e,\mathcal C<B^K$ and $L>3K+2$ give
$2e\mathcal C^2<q$.

Lemma~\ref{lem:borrow-high-window} applies with $Z=4$.
Its coefficient plateau is $J_{\mathcal C}$: the digit of $S_3$
at $L$ is $B-J_{\mathcal C}-1$ or $B-J_{\mathcal C}$, and the
digits at $L+1,\ldots,L+K$ are $B-J_{\mathcal C}$.
The high third-mask part is $X=(B-4)a_{\mathrm{al}}$, of degree at most $K$.
No overlap bounds every digit of $X$ by $J_{\mathcal C}$.
For its digit polynomial $F_X$, evaluation at four gives
\[
 B-4\mid F_X(4),\qquad
 0\le F_X(4)\le
 J_{\mathcal C}\frac{4^{K+1}-1}{3}<\frac B{12}<B-4.
\]
The second bound on $H$ proves the strict inequality. Hence $X=0$
and $a_{\mathrm{al}}=0$: both codes are canonical. Zero digits of $e_0$ have caused
no change, since the evaluation lemma requires digits in $[0,4)$ and
the high-window argument requires $e_0>0$, not positivity of each digit.

Below $K$, the raw coefficients of $S_3$ are those of
$2\mathcal D\mathcal C^2$. The non-strict coefficient bound
$|\mathcal D_j|\le2$ is sufficient because
\[
 \bigl|[B^j]\,2\mathcal D\mathcal C^2\bigr|
 \le4\mathcal C(1)^2=J_{\mathcal C}<B/4.
\]
All incoming carries in this range are zero or minus one.
The mask $B-4$ permits just the digits $0,1,2,3$; a negative raw
value plus its incoming carry normalizes to a digit greater than three.
The special target has raw coefficient $2\delta^2$, forcing
$\delta\in\{0,1\}$. If $\delta=0$, the guard has raw coefficient
$-4x^2$ and fails. Thus $\delta=1$.
Each ordinary target now has raw coefficient $4F_j+2$. With either
possible incoming carry, it is permitted exactly when $F_j=0$.
The original system therefore holds at the actual queried input $x$.

Conversely, extend an original solution by its gate values, take
$\delta=1$ and $u=x^2$, and choose a power-of-two $b$ larger than all
coordinates. Use the canonical codes. The low coordinate tests vanish
by support separation, and every ordinary or special target has raw
coefficient $2$, hence normalized digit $1$ or $2$.
The inequalities $e_0,\ell_0<q$ give $S_2<q^2$ and
$\Omega>2\lambda$. Since $2K<L$, one has
$4\mathcal C^2<Bq$, which gives $S_3>0$.
All required positive slack witnesses therefore exist.
The positive Pell witnesses from Section~\ref{sec:common-witness}
complete the final system. This proves both directions.

\begin{theorem}\label{thm:best101}
For every recursively enumerable subset of the positive integers,
there is an admissible fixed index $(V,H,L)$ such that its membership
relation is represented by a system with 34 positive unknowns and
22 equations admitting a straight-line certificate of 101 operations:
56 multiplications and 45 additions. The supplied parameters are
$(x,V,H,L)$, with three fixed index parameters. Fixed numerals are free.
\end{theorem}
\begin{proof}
The preceding construction proves the representation and both
positive-domain implications. In the 102-operation certificate,
replace $Z$ by $4$, replace the old exponent numeral by $L$, and
replace all uses of the separately computed $B-4$ by $\theta$.
The already present equation $\theta+4=B$ justifies that substitution.
It removes one subtraction, with no added arithmetic, unknown, or
equation. The exact primitive and polynomial-residual check is recorded
in \texttt{round23\_1980\_fixed\_four\_certificate.py} and its JSON
receipt. Thus the count is $56+45=101$; no optimality is asserted.
\end{proof}
